Simulators of large-language-model serving systems and RTL implementations of
RDMA endpoints both need the same thing and neither can buy it: a public,
mechanism-level account of how a commodity RoCEv2 NIC actually behaves. Vendor
silicon ships no architectural contract, published NIC designs are either
research prototypes or closed IP, and the constants that decide whether a
tensor-parallel all-reduce is bandwidth bound or loss bound are nowhere in a
datasheet.

This paper reconstructs one such device, the NVIDIA (Mellanox) ConnectX-5 Ex at
100\,GbE, from public information alone: the InfiniBand Architecture
Specification and its RoCEv2 annex, the vendor programmer's reference manual,
the open-source \texttt{mlx5} kernel driver and \texttt{rdma-core} provider,
the published literature, and more than 700 pre-registered black-box
measurements on a four-node cluster the authors are entitled to use. We
describe the campaign method, in which every phase freezes its expectations
before the run so that a miss is evidence rather than an embarrassment, and we
report what it found: a fixed-offset initiation cost of $4.48$\uS{} against a
$97.1$\Gbps{} goodput asymptote; queue depth dominating message size by a
factor of $5.9$ at 8\,KiB; a saturated single queue pair whose apparent
78\Gbps{} ``cap'' is a loss and retransmission equilibrium rather than a
hardware limit; a drain window above which a paced burst becomes exactly
lossless; a 195\Gbps{} full-duplex aggregate and a 197\Gbps{} internal
processing budget in which wire traffic outranks loopback; ECT(0) stamped
unconditionally on every transmitted RoCEv2 packet; and, on a fabric whose
switch was measured never to mark a single packet in 670 million, congestion
notifications generated by the receiving NIC from its own ingress congestion,
with reaction dynamics two to twenty times slower than the DCQCN defaults
suggest.

From those mechanisms we derive the paper's main artifact: a nineteen-block
architecture for an FPGA implementation on an AMD Alveo U55C, expandable to an
ASIC, carrying an AXI4 data path for host and memory traffic and a separate
AXI4-Lite monitoring network on chip that gives every block a slot in one
versioned control and status register map mirroring the observable state of the
real NIC. Each block records the measured behaviour it must reproduce, the
evidence class of every parameter it carries, and the block of the golden C
model it corresponds to. A nineteen-row anomaly table serves as the
verification contract: the golden model is driven through DPI-C from a UVM
environment, and a block is not validated until the anomaly registered against
it is reproduced inside its registered band.
